\documentclass[11pt]{article}
\usepackage[margin=25mm]{geometry}
\usepackage{amsmath,booktabs,siunitx}
\usepackage[hidelinks]{hyperref}
\usepackage[nameinlink,noabbrev]{cleveref}
\hypersetup{pdftitle={A guided reference document},pdfauthor={Course reader}}
\begin{document}
\tableofcontents
\section{Question}\label{sec:question}
How can a reader find an equation, a table, and the source of a command?
The details are in \cref{sec:method}; the arithmetic is in \cref{eq:average}.
The LaTeX Project provides official documentation \cite{latex-project}.

\section{Method}\label{sec:method}
Let $t_1$ and $t_2$ be the two task times in minutes.
The symbol $\bar t$ denotes their arithmetic mean:
\begin{equation}\label{eq:average}
\bar t=\frac{t_1+t_2}{2}.
\end{equation}
The values in \cref{tab:times} are invented teaching data.
They demonstrate alignment and references.
\begin{table}[htbp]
\centering
\caption{Invented task times}\label{tab:times}
\begin{tabular}{lS[table-format=2.2]}
\toprule
Task & {Time (min)} \\
\midrule
Draft & 9.50 \\
Review & 12.25 \\
\bottomrule
\end{tabular}
\end{table}
\Cref{eq:average} is on page~\pageref{eq:average}.
Return to \cref{sec:question} or read
\href{https://www.latex-project.org/help/documentation/}{the official documentation}.

\begin{thebibliography}{9}
\raggedright
\bibitem{latex-project}
The LaTeX Project.
\emph{Documentation}.
\url{https://www.latex-project.org/help/documentation/}.
\end{thebibliography}
\end{document}